\documentclass[11pt,a4paper]{moderncv}
\moderncvstyle{banking}
\moderncvcolor{blue}
\usepackage[scale=0.75]{geometry}
\usepackage{hyperref}
\usepackage{fontawesome5}


% -----------------------------
% Datos personales
% -----------------------------
\name{Matihus}{Molina Larios}
\address{Lima, Per\'u}{}
\email{matihus.molina@unmsm.edu.pe}
\social[linkedin]{maticus}
\social[github]{Maticus-7}

% -----------------------------
\begin{document} 
\makecvtitle

% -----------------------------
% Perfil profesional
% -----------------------------
\section{Perfil profesional}
\cvitem{}{ Soy estudiante de último año de matemáticas en la Universidad Nacional de San Marcos, con interés en la teoría ergódica, las aplicaciones matemáticas a la salud y las matemáticas formalmente verificadas.
	Recientemente, he desarrollado un interés particular en la formalización de las matemáticas mediante LEAN4. También participo en la preparación de artículos de investigación relacionados con el aprendizaje automático aplicado al cáncer de pulmón y un estudio comparativo de métodos numéricos para ecuaciones diferenciales de tipo Lane-Emden.
	Mi objetivo actual es fortalecer mi formación en matemáticas puras y aplicadas y contribuir al desarrollo de la comunidad académica en Perú. Tras graduarme, aspiro a realizar estudios de posgrado en el extranjero, explorando diferentes áreas de las matemáticas.}

% -----------------------------
% Educación
% -----------------------------
\section{Educaci\'on}
\cventry{03/2022 -- 12/2026 (estimado)}{Estudiante de Bachillerato en Ciencias Matem\'aticas}{Universidad Nacional Mayor de San Marcos}{Lima, Per\'u}{}{Promedio general (CGPA): 15.19/20}
\cventry{2018 -- 2021}{Egresado del Programa de Ingl\'es Integral}{Centro Cultural LIDEMPER\'U}{Lima, Per\'u}{}{}

% -----------------------------
% Idiomas
% -----------------------------
\section{Idiomas}
\cvitem{Espa\~nol}{Nativo}
\cvitem{Ingl\'es}{Intermedio B1}
\cvitem{Portugu\'es}{B\'asico A2}

% -----------------------------
% Proyectos académicos
% -----------------------------
\section{Proyectos académicos}
\cventry{01/2025 -- 03/2025}{Proyecto: ``Un primer encuentro con los sistemas dinámicos caóticos''}{PICV-UNMSM}{}{}{
\begin{itemize}
  \item Docente asesor: Dr. Jorge L. Crisóstomo
  \item Estudio de la conjugación topológica entre $Q_c$ y el shift de Bernoulli.
  \item Aplicación de conceptos de ecuaciones diferenciales, análisis funcional y topología.
  \item \href{http://bit.ly/3Hj8NvJ}{Click para verlo}
\end{itemize}}

\cventry{02/2025 -- 03/2025}{Predicción de Depresión en Estudiantes}{Curso: Matemáticas para ML y DS (UNMSM)}{}{}{
\begin{itemize}
  \item Aplicación de probabilidad, estadística y técnicas de machine learning.
  \item Modelado supervisado con Python y scikit-learn.
  \item \href{https://github.com/Maticus-7/Mental-Healthy/tree/2bc0dac3f89d9187dfa37b6cd963d69256551167/Depression/University}{Click para verlo}
\end{itemize}}

\cventry{10/2025 --12/2025}{Numerical solution for differential equations of Lane-Emden type}{Modelamiento de Cuerpos Estelares}{}{}{
\begin{itemize}
  \item Análisis de los métodos de descomposición y sus diferencias. 
  \item Modelamiento númerico de los métodos y su eficacia. 
  \item \href{https://github.com/Maticus-7/lane-emden}{Click para verlo}
\end{itemize}}

% -----------------------------
% Eventos académicos
% -----------------------------
\section{Eventos acad\'emicos}
%\cventry{Agosto de 2025}{PUCP-Lima}{XXXVIII COLOQUIO DE LA SOCIEDAD MATEMÁTICA PERUANA}{}{}{Asistencia a ponencias sobre topolog\'ia, an\'alisis funcional, geometr\'ia diferencial, sistemas din\'amicos, entre otros temas de investigaci\'on actual.
%Fue una gran oportunidad para conocer m\'as sobre el trabajo de matem\'aticos nacionales e internacionales, as\'i como entablar conversaciones con investigadores y profesores destacados.}

\cventry{Octubre de 2025}{Concepción,Chile}{E'(UBB)=2025}{}{}{
Como primera experiencia internacional fue inolvidable,me encantó la ciudad y el conocer otra cultura.Me parecío muy interesante la idea de hacer la maestría allí.}

\cventry{Diciembre de 2025}{Huaraz,Perú}{VII EICM}{}{}{
Es mi primera experiencia presentando un proyecto de investigación en un evento de esta categoría. El proyecto en cuestión es el de Modelamiento de Cuerpo Estelares. }

\cventry{Febrary 2026}{UNT -- Trujillo}{Summer School and Workshop on Optimization and Operator Learning}{}{}{Participé en este importante evento organizado por el SaC$^3$, poniendo a prueba mis conocimientos en teoría de optimización y programación, destacando en la pre-escuela.}

% -----------------------------
% Áreas de interés
% -----------------------------
\section{\'Areas de inter\'es}
\cvitem{Matem\'atica Pura}{
\begin{itemize}
    \item An\'alisis Funcional.
    \item Teor\'ia Erg\'odica
    \item Fundamentos Matem\'aticos de la Inteligencia Artificial.
    \item Matem\'aticas Formalmente Verificadas. 
    \item Aplicación matemática a la salud humana.
\end{itemize}}
  
\cvitem{Tecnolog\'ias}{
\begin{itemize}
\item Lenguajes: Python, R, MATLAB, LEAN4
\item  Librer\'ias:  NumPy,SciPy,Pandas, Scikit-Learn, Matplotlib,PyTorch.
\item Herramientas: Git, \LaTeX, Power BI.
\end{itemize}}

	% ------------------------------
% Logros
% ------------------------------
\section{Premios y Distinciones}
\cvitem{}{BECA PERMANENCIA 2025 ,PRONABEC}
\cvitem{}{Reconocimiento de mi facultad por mi EXCELENCÍA ACADÉMICA}  


% -----------------------------
% Actividades extracurriculares
% -----------------------------

\section{Actividades extracurriculares}
\cvitem{}{Colaborador activo del IEEE Computational Intelligence Society – UNMSM}
\cvitem{}{Exmiembro de la agrupación cultural y artística de salsa SALSUNI – San Marcos}
\cvitem{}{Exmiembro del equipo universitario de judo – Universidad Nacional Mayor de San Marcos}
\cvitem{}{Miembro activo del BILDHEIT}

% -----------------------------
% Certificados
% -----------------------------
% Sección: Certificados
\section{Certificados} 
\cvitem{}{ \href{https://maticus-7.github.io/Certificates/}{Los certificados de cursos y logros académicos están disponibles en el siguiente repositorio}.}


\end{document}



